\documentclass[12pt]{article}

\usepackage[margin=1.12in]{geometry}
\usepackage{amsmath,amssymb,amsthm,mathtools}
\usepackage{booktabs}
\usepackage{hyperref}

\newtheorem{theorem}{Theorem}
\newtheorem{proposition}{Proposition}
\theoremstyle{remark}
\newtheorem{remark}{Remark}

\newcommand{\A}{\mathcal A}
\newcommand{\B}{\mathcal B}
\newcommand{\Hcal}{\mathcal H}
\newcommand{\W}{\mathcal W}
\newcommand{\barotimes}{\mathbin{\bar\otimes}}
\newcommand{\id}{\operatorname{id}}

\setlength{\emergencystretch}{3em}
\hypersetup{
  pdftitle={Generated-Algebra Correlations Reconstruct Tensor Factors Only After Split or Kernel Separation},
  pdfsubject={split inclusions, tensor products, finite centers, AQFT},
  pdfkeywords={split property, tensor product, von Neumann algebra, multiplication kernel, finite center}
}

\title{Generated-Algebra Correlations Reconstruct\\
Tensor Factors Only After Split or Kernel Separation}
\author{}
\date{Live review note of June 20, 2026}

\begin{document}
\maketitle

\begin{abstract}
Split inclusions and W*-tensor-product independence are standard positive
hypotheses in algebraic quantum field theory: when the multiplication map
\(\mu:\A\bar\otimes\B\to\A\vee\B\) is a normal isomorphism, the generated
algebra determines the spatial tensor-product coordinate.  This note isolates
the converse boundary.  If the supplied data are only generated-algebra
correlations \(\gamma(\mu(x))\), then a candidate tensor-product element
\(x\in\A\bar\otimes\B\) is reconstructed exactly modulo the common tested
multiplication kernel.  In particular, a shared nontrivial central projection
\(p\in\A\cap\B\) gives \(0\ne p\otimes(1-p)\in\ker\mu\).  Thus
generated-algebra tomography is not, by itself, a split/tensor-factor theorem.
\end{abstract}

\paragraph{Status.}
This is non-frozen external-review metadata for Team 1.  It is intended as a
small claim that an operator-algebra or AQFT reviewer can classify as known,
false, missing-hypothesis, or useful.  It does not modify any frozen Team 1
manuscript, Lean archive, public mirror manifest, or SHA-256 package.

\section{Primary sources}

The boundary below was checked against Doplicher--Longo split inclusions
\cite{DoplicherLongo}, Buchholz--Wichmann nuclearity and causal independence
\cite{BuchholzWichmann}, and Werner local preparability
\cite{WernerSplit}.  These works are used here as positive routes to split
independence, not as claims that split theory is absent from physics.

\section{Generated algebra versus spatial tensor product}

Let \(\A,\B\subset B(\Hcal)\) be commuting von Neumann algebras, and let
\(\A\vee\B\) be the von Neumann algebra they generate.  Suppose multiplication
extends to a normal representation
\[
  \mu:\A\barotimes\B\longrightarrow \A\vee\B,
  \qquad
  \mu(a\otimes b)=ab .
\]
Let \(\Gamma\subset(\A\vee\B)_*\) be the normal functionals used as supplied
generated-algebra correlations:
\[
  D_\Gamma(x)=\bigl(\gamma(\mu(x))\bigr)_{\gamma\in\Gamma},
  \qquad x\in\A\barotimes\B .
\]

\begin{theorem}[multiplication-kernel criterion]
\label{thm:multiplication-kernel}
The data \(D_\Gamma\) reconstruct the spatial tensor-product element \(x\)
if and only if
\[
  \bigcap_{\gamma\in\Gamma}\ker(\gamma\circ\mu)=\{0\}.
  \tag{1}
\]
If \(\Gamma\) separates points of \(\A\vee\B\), then (1) is equivalent to
\[
  \ker\mu=\{0\}.
  \tag{2}
\]
Thus full generated-algebra tomography reconstructs the unquotiented spatial
tensor-product coordinate exactly when the multiplication representation is
injective.
\end{theorem}

\begin{proof}
Two candidates \(x,y\in\A\barotimes\B\) have the same supplied data exactly
when
\[
  \gamma(\mu(x-y))=0\qquad(\gamma\in\Gamma),
\]
that is, when \(x-y\) lies in the common kernel in (1).  Therefore the data
separate all candidates precisely when that kernel is zero.  If \(\Gamma\)
separates \(\A\vee\B\), then \(\gamma(\mu(z))=0\) for all \(\gamma\in\Gamma\)
is equivalent to \(\mu(z)=0\), giving (2).
\end{proof}

\begin{proposition}[shared-center witness]
\label{prop:shared-center}
If \(\A\cap\B\) contains a projection \(0\ne p\ne1\), then
\[
  0\ne p\otimes(1-p)\in\A\barotimes\B,
  \qquad
  \mu\bigl(p\otimes(1-p)\bigr)=0 .
\]
Consequently no data package that factors only through \(\A\vee\B\) can
distinguish the zero tensor-product element from this nonzero tensor-kernel
direction.
\end{proposition}

\begin{proof}
The projections \(p\) and \(1-p\) are nonzero in the two tensor factors, so
\(p\otimes(1-p)\) is a nonzero positive element of the spatial tensor product.
Because \(p\) is a common projection, multiplication sends it to
\(p(1-p)=0\).  Any datum that factors through \(\mu\) assigns the same value to
this element and to \(0\).
\end{proof}

\section{Positive exits}

\begin{center}
\begin{tabular}{p{0.34\linewidth}p{0.51\linewidth}}
\toprule
Exit premise & Effect \\
\midrule
Split inclusion or type-I interpolation & Supplies the standard
W*-tensor-product independence structure for the bipartition. \\
\addlinespace
Normal isomorphism \(\A\barotimes\B\simeq\A\vee\B\) & Makes \(\mu\) injective;
the inverse of \(\mu\) reconstructs the tensor coordinate. \\
\addlinespace
Kernel-separating tests & Adds data that distinguish directions in
\(\ker\mu\). \\
\addlinespace
Quotient, factor, or inadmissibility axiom & Removes the kernel directions
from the physical target instead of reconstructing them. \\
\bottomrule
\end{tabular}
\end{center}

The exact external question is:

\begin{quote}
Does the target AQFT setup assume split/W*-tensor-product independence, type-I
interpolation, factor-valuedness, or kernel-separating tests for the
bipartition?  If not, generated-algebra correlations reconstruct only the
quotient of the candidate spatial tensor product by \(\ker\mu\).
\end{quote}

If the split hypothesis is supplied, it is a positive exit from the Team 1
finite-center obstruction.  If it is absent, the Team 1 row is only the
standard multiplication-kernel boundary, not a new split theorem.

\begin{thebibliography}{9}

\bibitem{DoplicherLongo}
S. Doplicher and R. Longo,
\emph{Standard and split inclusions of von Neumann algebras},
Invent. Math. 75 (1984), 493--536,
\href{https://link.springer.com/article/10.1007/BF01388641}{doi:10.1007/BF01388641}.

\bibitem{BuchholzWichmann}
D. Buchholz and E. H. Wichmann,
\emph{Causal independence and the energy-level density of states in local quantum field theory},
Commun. Math. Phys. 106 (1986), 321--344,
\href{https://link.springer.com/article/10.1007/BF01454978}{doi:10.1007/BF01454978}.

\bibitem{WernerSplit}
R. F. Werner,
\emph{Local preparability of states and the split property in quantum field theory},
Lett. Math. Phys. 13 (1987), 325--329,
\href{https://link.springer.com/article/10.1007/BF00401161}{doi:10.1007/BF00401161}.

\end{thebibliography}

\end{document}
